\documentclass[a4paper, zihao=-4]{ctexart}

% ==================== 宏包配置 ====================
\usepackage[utf8]{inputenc}
\usepackage[T1]{fontenc}
\usepackage{geometry}
\geometry{left=2.5cm, right=2.5cm, top=2.5cm, bottom=2.5cm} % 标准页边距

\ctexset{
  section/format += \raggedright,
  subsection/format += \raggedright
}

% 行距设置
\usepackage{setspace}
\setstretch{1.5} % 设置1.5倍行距

% 代码高亮设置
\usepackage{xcolor}
\usepackage{listings}
\definecolor{codegreen}{rgb}{0,0.6,0}
\definecolor{codegray}{rgb}{0.5,0.5,0.5}
\definecolor{codepurple}{rgb}{0.58,0,0.82}
\definecolor{backcolour}{rgb}{0.96,0.96,0.96}

\lstset{
    backgroundcolor=\color{backcolour},   
    commentstyle=\color{codegreen},
    keywordstyle=\color{magenta},
    numberstyle=\tiny\color{codegray},
    stringstyle=\color{codepurple},
    basicstyle=\ttfamily\footnotesize, % 代码字体稍小以适应版面
    breakatwhitespace=false,         
    breaklines=true,                 
    captionpos=b,                    
    keepspaces=true,                 
    numbers=left,                    
    numbersep=5pt,                  
    showspaces=false,                
    showstringspaces=false,
    showtabs=false,                  
    tabsize=4,
    frame=single,
    language=Python,
    extendedchars=true,
    escapeinside=``
}

% 绘图与图片支持
\usepackage{graphicx}
\usepackage{tikz}
\usetikzlibrary{shapes, arrows, positioning, fit, calc}
\usepackage{float} % 用于固定图片位置

% 链接设置
\usepackage{hyperref}
\hypersetup{
    colorlinks=true,
    linkcolor=black,
    filecolor=black,      
    urlcolor=blue,
    pdftitle={Kea2性质检查模块技术报告},
}

% 标题本地化
\renewcommand{\figurename}{图}
\renewcommand{\tablename}{表}

% 简单的Verbatim环境优化
\usepackage{fancyvrb}
\usepackage{enumitem}
\setlist[enumerate,itemize]{noitemsep, topsep=0pt}

\newcommand{\bff}[1]{{\heiti #1}}

% ==================== 文档开始 ====================
\title{\textbf{Kea2性质检查模块技术报告}}
\author{}
\date{}

\begin{document}

\maketitle

\section{技术原理复盘}

\subsection{模块目标与解决的问题}

Kea2的性质检查模块旨在解决Android应用非崩溃功能bug（占APP总bug的65.4\%）难以自动化检测的问题。该模块通过\bff{基于性质的测试（Property-based Testing, PBT）}，在APP动态执行过程中自动验证用户指定的程序性质，精准识别功能异常。

\bff{核心解决的问题}：
\begin{itemize}
    \item \bff{测试预言缺失}：传统自动化GUI测试缺乏有效机制判断功能是否正确。
    \item \bff{手动测试成本高}：功能bug检测依赖人工验证，效率低下。
    \item \bff{覆盖不足}：现有工具仅能覆盖21\%的功能bug，且输入质量低。
\end{itemize}

\subsection{输入与输出}

\bff{输入}：
\begin{enumerate}
    \item \bff{性质规约}：用户通过PDL（Property Description Language，基于Python的领域特定语言）编写的性质定义，包含：
    \begin{itemize}
        \item \bff{前置条件（Precondition, P）}：使用 \texttt{@precondition} 装饰器定义，指定性质可执行时的GUI状态。
        \item \bff{交互场景（Interaction, I）}：性质测试方法体中的操作序列（如点击按钮、输入文本、旋转屏幕等）。
        \item \bff{后置条件/断言（Postcondition/Assertion, Q）}：在测试方法中使用 \texttt{assert} 语句定义，验证交互后的程序状态。
    \end{itemize}
    \item \bff{待测APK文件}：目标Android应用的安装包。
    \item \bff{探索策略配置}：随机探索或主路径引导探索的参数设置。
\end{enumerate}

\bff{输出}：
\begin{enumerate}
    \item \bff{性质检查结果}：通过（pass）/违反（fail/error）。
    \item \bff{Bug报告}：违反性质时生成HTML格式的详细报告，包含：
    \begin{itemize}
        \item 触发bug的GUI测试用例（事件轨迹）。
        \item 程序状态快照（截图）。
        \item 性质执行统计信息。
        \item 堆栈轨迹（traceback）。
    \end{itemize}
\end{enumerate}

\subsection{基本流程与代码架构}

\subsubsection{基本流程（三阶段）}

\bff{阶段(a)：获取用户定义的性质作为输入}

\begin{lstlisting}
# keaUtils.py:419-422
def run(self, test):
    self.allProperties = dict()
    self.collectAllProperties(test)  # 收集所有标记了@precondition的测试方法
\end{lstlisting}

\begin{itemize}
    \item \bff{入口}：\texttt{KeaTestRunner.run()} 方法。
    \item \bff{机制}：通过反射遍历所有 \texttt{unittest.TestCase}，查找标记了 \texttt{PRECONDITIONS\_MARKER}（即 \texttt{@precondition} 装饰器）的测试方法。
    \item \bff{存储}：将性质存储到 \texttt{self.allProperties} 字典中，键为方法名，值为 \texttt{TestCase} 实例。
\end{itemize}

\bff{阶段(b)：当满足前置条件时，在动态执行过程中暂停，将控制权转移以执行性质检查}

这是模块的核心机制，具体实现如下：

\begin{lstlisting}
# keaUtils.py:474-544
while self.stepsCount < self.options.maxStep:
    # 1. 获取当前GUI状态（通过Fastbot发送Monkey事件或直接dump hierarchy）
    if fb.executed_prop:
        # 如果刚执行完性质，直接获取当前状态，不发送新事件
        xml_raw = fb.dumpHierarchy()
    else:
        # 正常探索：发送Monkey事件并获取GUI状态
        self.stepsCount += 1
        xml_raw = fb.stepMonkey(self._monkeyStepInfo)
    
    # 2. 检查所有性质的前置条件（暂停探索，进行状态检查）
    propsSatisfiedPrecond = self.getValidProperties(xml_raw, result)
    
    # 3. 如果没有满足前置条件的性质，继续探索（控制权返回给探索循环）
    if len(propsSatisfiedPrecond) == 0:
        continue
    
    # 4. 有满足前置条件的性质时，暂停探索，转移控制权执行性质检查
    # ... 概率过滤和选择性质 ...
    execPropName = random.choice(propsNameFilteredByP)
    test = propsSatisfiedPrecond[execPropName]
    
    # 5. 执行性质测试方法（控制权转移到用户定义的性质代码）
    test(result)  # 执行交互场景I和断言Q
    
    # 6. 性质执行完成后，标记已执行，下次循环直接获取状态
    fb.executed_prop = True
\end{lstlisting}

\bff{暂停机制的关键点}：
\begin{itemize}
    \item \bff{暂停探索}：当检测到满足前置条件的性质时，不再继续发送Monkey事件（\texttt{stepsCount} 不增加）。
    \item \bff{控制权转移}：从Fastbot的探索循环转移到性质测试方法的执行（\texttt{test(result)} 调用）。
    \item \bff{状态保持}：通过 \texttt{fb.executed\_prop} 标志，确保性质执行后能获取当前GUI状态，而不是立即发送新事件。
\end{itemize}

\bff{阶段(c)：根据用户定义的"前置条件-后置条件(断言)"进行性质检查，并将结果记录为输出}

\begin{lstlisting}
# keaUtils.py:536-544
try:
    test(result)  # 执行性质测试方法
    # 方法体内包含：
    # - 交互场景I：self.d(text="Settings").click() 等操作
    # - 后置条件Q：assert self.d(...).exists() 等断言
finally:
    result.printError(test)

# 记录执行结果
result.updateExectedInfo()  # 更新状态为pass/fail/error
fb.logScript(result.lastExecutedInfo)  # 记录到Fastbot日志
result.flushResult()  # 持久化到JSON文件
\end{lstlisting}

\begin{itemize}
    \item \bff{断言评估}：性质测试方法中的 \texttt{assert} 语句如果失败，会抛出 \texttt{AssertionError}。
    \item \bff{异常捕获}：\texttt{unittest.TestCase} 的 \texttt{\_\_call\_\_} 方法捕获异常，调用 \texttt{result.addFailure()} 或 \texttt{result.addError()}。
    \item \bff{结果记录}：执行信息（开始步骤数、状态、堆栈轨迹）记录到 \texttt{PropertyExecutionInfo}，并持久化到文件。
\end{itemize}

\subsubsection{代码架构}

\bff{主要类/函数分工}：

\begin{enumerate}
    \item \bff{PDL解析器}（\texttt{keaUtils.py} 中的装饰器系统）：
    \begin{itemize}
        \item \texttt{@precondition}（\texttt{keaUtils.py:54-65}）：解析前置条件，支持多个前置条件（逻辑AND）。
        \item \texttt{@prob}（\texttt{keaUtils.py:68-81}）：定义性质执行概率（0-1之间）。
        \item \texttt{@max\_tries}（\texttt{keaUtils.py:84-97}）：限制性质最大执行次数。
        \item \texttt{collectAllProperties()}（\texttt{keaUtils.py:631-672}）：通过反射收集所有标记了 \texttt{@precondition} 的测试方法。
    \end{itemize}
    
    \item \bff{探索策略模块}（\texttt{fastbotManager.py} + \texttt{keaUtils.py}）：
    \begin{itemize}
        \item \texttt{FastbotManager}：管理Fastbot服务，发送Monkey事件，获取GUI状态。
        \item \texttt{KeaTestRunner.run()}：主循环，协调探索和性质检查。
    \end{itemize}
    
    \item \bff{状态收集器}（\texttt{u2Driver.py}）：
    \begin{itemize}
        \item \texttt{U2StaticChecker}：基于XML hierarchy进行静态检查（用于前置条件评估）。
        \item \texttt{U2ScriptDriver}：可执行实际操作的驱动（用于性质执行时的交互）。
    \end{itemize}
    
    \item \bff{断言评估器}（\texttt{keaUtils.py} 中的 \texttt{KeaTestRunner}）：
    \begin{itemize}
        \item \texttt{JsonResult.addFailure/addError()}：捕获断言失败和异常，记录执行信息。
    \end{itemize}
    
    \item \bff{结果记录与报告生成}（\texttt{keaUtils.py} + \texttt{report/bug\_report\_generator.py}）：
    \begin{itemize}
        \item \texttt{JsonResult}：实时记录性质执行统计。
        \item \texttt{BugReportGenerator}：生成HTML格式的bug报告。
    \end{itemize}
\end{enumerate}

\section{源码阅读与逻辑拆解}

\subsection{完整逻辑链：从输入到输出}

\subsubsection{输入阶段：性质定义与收集}

\bff{入口函数}：\texttt{kea\_launcher.py:run()} $\rightarrow$ \texttt{KeaTestRunner.run()}

\bff{性质收集流程}（\texttt{keaUtils.py:631-672}）：

\begin{lstlisting}
def collectAllProperties(self, test: TestSuite):
    def iter_tests(suite):
        for test in suite:
            if isinstance(test, TestSuite):
                yield from iter_tests(test)
            else:
                yield test
    
    # 遍历所有TestCase，查找标记了@precondition的测试方法
    for t in iter_tests(test):
        testMethodName = t._testMethodName
        testMethod = getattr(t, testMethodName)
        
        # 检查是否有PRECONDITIONS_MARKER标记（即@precondition装饰器）
        if hasattr(testMethod, PRECONDITIONS_MARKER):
            # 移除setUp/tearDown钩子（PBT不需要）
            remove_setUp(t)
            remove_tearDown(t)
            # 保存到allProperties字典
            self.allProperties[testMethodName] = t
\end{lstlisting}

\bff{性质定义示例}（\texttt{properties/Omninotes\_Sample.py}）：

\begin{lstlisting}
@prob(0.7)  # 执行概率70%
@precondition(
    lambda self: self.d(text="Omni Notes Alpha").exists
    and self.d(text="Settings").exists
)
def test_goToPrivacy(self):
    """交互场景：点击Settings -> Privacy"""
    self.d(text="Settings").click()
    sleep(0.5)
    self.d(text="Privacy").click()

@max_tries(1)  # 最多执行1次
@precondition(
    lambda self: self.d(resourceId="it.feio.android.omninotes.alpha:id/search_src_text").exists
)
def test_rotation(self):
    """交互场景：旋转设备，断言：搜索框应仍然存在"""
    self.d.set_orientation("l")
    sleep(2)
    self.d.set_orientation("n")
    sleep(2)
    # 后置条件/断言Q
    assert self.d(resourceId="it.feio.android.omninotes.alpha:id/search_src_text").exists()
\end{lstlisting}

\subsubsection{处理阶段：探索循环与前置条件检查}

\bff{主探索循环}（\texttt{keaUtils.py:474-544}）：

\begin{lstlisting}
while self.stepsCount < self.options.maxStep:
    # === 步骤1：获取GUI状态 ===
    if fb.executed_prop:
        # 刚执行完性质，直接获取当前状态（不发送新事件）
        fb.executed_prop = False
        xml_raw = fb.dumpHierarchy()  # 通过HTTP请求获取XML hierarchy
    else:
        # 正常探索：发送Monkey事件
        self.stepsCount += 1
        xml_raw = fb.stepMonkey(self._monkeyStepInfo)  # 发送事件并获取GUI状态
    
    # === 步骤2：检查所有性质的前置条件（暂停探索） ===
    propsSatisfiedPrecond = self.getValidProperties(xml_raw, result)
    
    # === 步骤3：如果没有满足前置条件的性质，继续探索 ===
    if len(propsSatisfiedPrecond) == 0:
        continue
    
    # === 步骤4：概率过滤 ===
    p = random.random()
    propsNameFilteredByP = []
    for propName, test in propsSatisfiedPrecond.items():
        result.addPrecondSatisfied(test)  # 记录前置条件满足次数
        if getattr(test, PROB_MARKER, 1) >= p:
            propsNameFilteredByP.append(propName)
    
    if len(propsNameFilteredByP) == 0:
        continue
    
    # === 步骤5：随机选择一个性质执行（控制权转移） ===
    execPropName = random.choice(propsNameFilteredByP)
    test = propsSatisfiedPrecond[execPropName]
    
    # 依赖注入：将可执行操作的driver注入到test实例
    self.scriptDriver = U2Driver.getScriptDriver(mode="proxy")
    setattr(test, self.options.driverName, self.scriptDriver)
    
    # === 步骤6：执行性质测试方法 ===
    result.addExcuted(test, self.stepsCount)  # 记录执行开始
    fb.logScript(result.lastExecutedInfo)
    try:
        test(result)  # 执行交互场景I和断言Q
    finally:
        result.printError(test)
    
    # === 步骤7：记录结果 ===
    result.updateExectedInfo()  # 更新状态（pass/fail/error）
    fb.logScript(result.lastExecutedInfo)
    fb.executed_prop = True  # 标记已执行性质
    result.flushResult()  # 持久化结果到JSON文件
\end{lstlisting}

\bff{前置条件检查实现}（\texttt{keaUtils.py:589-629}）：

\begin{lstlisting}
def getValidProperties(self, xml_raw: str, result: JsonResult) -> PropertyStore:
    # 创建静态检查器（基于XML hierarchy，不连接设备，性能高）
    staticCheckerDriver = U2Driver.getStaticChecker(hierarchy=xml_raw)
    
    validProps: PropertyStore = dict()
    for propName, test in self.allProperties.items():
        valid = True
        prop = getattr(test, propName)  # 获取测试方法对象
        
        # 检查所有前置条件（逻辑AND，所有precond都必须返回True）
        for precond in prop.preconds:  # preconds是@precondition装饰器添加的元组
            # 依赖注入：将静态检查器注入到test实例（用于precond lambda中访问self.d）
            setattr(test, self.options.driverName, staticCheckerDriver)
            try:
                # 执行前置条件lambda函数：precond(test)
                # 例如：lambda self: self.d(text="Settings").exists
                if not precond(test):
                    valid = False
                    break
            except u2.UiObjectNotFoundError:
                valid = False
                break
            except Exception as e:
                logger.error(f"Error when checking precond: {getFullPropName(test)}")
                valid = False
                break
        
        # 如果所有前置条件都满足，检查是否达到最大执行次数
        if valid:
            if result.getExcuted(test) >= getattr(prop, MAX_TRIES_MARKER, float("inf")):
                continue
            validProps[propName] = test
    
    staticCheckerDriver.clear_cache()  # 清理缓存
    return validProps
\end{lstlisting}

\bff{状态收集实现}（\texttt{u2Driver.py:515-543}）：

\begin{lstlisting}
class U2StaticChecker:
    def __init__(self):
        self.d = U2StaticDevice(U2ScriptDriver().getInstance())
    
    def setHierarchy(self, hierarchy: str):
        # 解析XML hierarchy为etree.Element
        self.d.xml = etree.fromstring(hierarchy.encode("utf-8"))
        # 过滤被覆盖的控件（hidden widget filter）
        _HindenWidgetFilter(self.d.xml)
    
    def getInstance(self, hierarchy: str=None):
        self.setHierarchy(hierarchy)
        return self.d  # 返回U2StaticDevice实例
\end{lstlisting}

\bff{控件存在性检查}（\texttt{u2Driver.py:192-197}）：

\begin{lstlisting}
@property
def exists(self):
    set_covered_to_deepest_node(self.selector)  # 设置covered属性
    xpath = self.selector_to_xpath(self.selector)  # 转换U2 Selector为XPath
    matched_widgets = self.session.xml.xpath(xpath)  # 在XML中查找
    return bool(matched_widgets)
\end{lstlisting}

\subsubsection{输出阶段：性质执行与断言评估}

\bff{性质执行机制}（\texttt{keaUtils.py:536-544}）：

\begin{lstlisting}
# 依赖注入：将脚本驱动（可执行操作的driver）注入到test实例
self.scriptDriver = U2Driver.getScriptDriver(mode="proxy")
setattr(test, self.options.driverName, self.scriptDriver)

# 记录执行开始
result.addExcuted(test, self.stepsCount)
fb.logScript(result.lastExecutedInfo)

try:
    # 执行性质测试方法（控制权转移到用户代码）
    test(result)  # 调用unittest.TestCase.__call__()
    # 方法体内执行：
    # 1. 交互场景I：self.d(text="Settings").click() 等操作
    # 2. 后置条件Q：assert self.d(...).exists() 等断言
finally:
    result.printError(test)  # 如果有错误，打印堆栈轨迹

# 更新执行信息（pass/fail/error）
result.updateExectedInfo()
fb.logScript(result.lastExecutedInfo)
fb.executed_prop = True  # 标记已执行性质，下次循环直接dump hierarchy
result.flushResult()  # 持久化结果到JSON文件
\end{lstlisting}

\bff{断言评估机制}：

\begin{itemize}
    \item 性质测试方法中的 \texttt{assert} 语句如果失败，会抛出 \texttt{AssertionError}。
    \item \texttt{unittest.TestCase} 的 \texttt{\_\_call\_\_} 方法会捕获异常，调用 \texttt{result.addFailure()}。
    \item \texttt{JsonResult.addFailure()}（\texttt{keaUtils.py:350-354}）记录失败状态和堆栈轨迹：
\end{itemize}

\begin{lstlisting}
def addFailure(self, test, err):
    super().addFailure(test, err)
    self.res[getFullPropName(test)].fail += 1
    self.lastExecutedInfo.state = "fail"
    self.lastExecutedInfo.tb = self._exc_info_to_string(err, test)
\end{lstlisting}

\bff{结果持久化}（\texttt{keaUtils.py:324-335}）：

\begin{lstlisting}
def flushResult(self):
    global RESFILE, PROP_EXEC_RESFILE
    # 持久化统计信息到result_*.json
    json_res = dict()
    for propName, propStatitic in self.res.items():
        json_res[propName] = asdict(propStatitic)
    with open(RESFILE, "w", encoding="utf-8") as fp:
        json.dump(json_res, fp, indent=4)
    
    # 持久化执行信息到property_exec_info_*.json（每行一个JSON对象）
    while self.executionInfoStore:
        execInfo = self.executionInfoStore.popleft()
        with open(PROP_EXEC_RESFILE, "a", encoding="utf-8") as fp:
            fp.write(f"{json.dumps(asdict(execInfo))}\n")
\end{lstlisting}

\subsection{模块执行流程图}

\begin{figure}[H]
\centering
\begin{tikzpicture}[node distance=1.5cm, auto,
    block/.style={rectangle, draw, fill=blue!10, text width=6em, text centered, rounded corners, minimum height=3em, font=\scriptsize},
    decision/.style={diamond, draw, fill=green!10, text width=5em, text centered, inner sep=0pt, font=\scriptsize},
    line/.style={draw, -latex},
    cloud/.style={draw, ellipse, fill=red!10, node distance=3cm, minimum height=2em}]

    % Nodes
    \node [block] (init) {KeaTestRunner\\run};
    \node [block, below of=init] (collect) {collectAllProperties};
    \node [block, below of=collect] (startFB) {FastbotManager\\start/init};
    \node [block, below of=startFB] (loop) {探索循环\\step < max?};
    
    \node [decision, below of=loop, node distance=2cm] (checkExec) {executed\_prop?};
    \node [block, left of=checkExec, node distance=3cm] (dump) {dumpHierarchy\\(获取状态)};
    \node [block, right of=checkExec, node distance=3cm] (step) {stepMonkey\\(发送事件)};
    
    \node [block, below of=checkExec, node distance=2.5cm] (checkPre) {getValidProperties\\(检查前置条件)};
    
    \node [decision, below of=checkPre, node distance=2cm] (hasValid) {满足前置条件?};
    
    \node [block, below of=hasValid, node distance=2cm] (execProp) {执行性质\\test(result)};
    \node [block, right of=execProp, node distance=3.5cm] (continue) {继续探索};
    
    \node [block, below of=execProp, node distance=2cm] (record) {记录结果\\flushResult};
    \node [block, below of=record] (setFlag) {set executed\_prop};

    % Edges
    \path [line] (init) -- (collect);
    \path [line] (collect) -- (startFB);
    \path [line] (startFB) -- (loop);
    
    \path [line] (loop) -- (checkExec);
    \path [line] (checkExec) -- node {True} (dump);
    \path [line] (checkExec) -- node {False} (step);
    
    \path [line] (dump) |- (checkPre);
    \path [line] (step) |- (checkPre);
    
    \path [line] (checkPre) -- (hasValid);
    \path [line] (hasValid) -- node {Yes} (execProp);
    \path [line] (hasValid) -| node [near start] {No} (continue);
    
    \path [line] (continue) |- (loop);
    \path [line] (execProp) -- (record);
    \path [line] (record) -- (setFlag);
    
    % Feedback loop
    \path [line] (setFlag) -- ++(-4,0) |- (loop);

\end{tikzpicture}
\caption{Kea2模块执行流程图}
\label{fig:flowchart}
\end{figure}

\bff{关键函数调用关系标注}：
\begin{itemize}
    \item \bff{性质收集}：\texttt{collectAllProperties()} $\rightarrow$ 反射遍历TestCase $\rightarrow$ 检查 \texttt{PRECONDITIONS\_MARKER}。
    \item \bff{状态获取}：\texttt{FastbotManager.stepMonkey()} $\rightarrow$ HTTP请求 $\rightarrow$ 返回XML hierarchy。
    \item \bff{前置条件检查}：\texttt{getValidProperties()} $\rightarrow$ \texttt{U2Driver.getStaticChecker()} $\rightarrow$ \texttt{U2StaticChecker.setHierarchy()} $\rightarrow$ \texttt{StaticU2UiObject.exists} $\rightarrow$ XPath查询。
    \item \bff{性质执行}：\texttt{test(result)} $\rightarrow$ \texttt{unittest.TestCase.\_\_call\_\_()} $\rightarrow$ 执行方法体 $\rightarrow$ 捕获异常 $\rightarrow$ \texttt{result.addFailure/addError()}。
    \item \bff{结果持久化}：\texttt{result.flushResult()} $\rightarrow$ 写入 \texttt{result\_*.json} 和 \texttt{property\_exec\_info\_*.json}。
\end{itemize}

\bff{数据流转路径标注}：
\begin{itemize}
    \item \bff{性质对象}：用户文件 $\rightarrow$ \texttt{allProperties} 字典 $\rightarrow$ \texttt{validProps} 字典 $\rightarrow$ \texttt{test} 实例 $\rightarrow$ \texttt{PropertyExecutionInfo} $\rightarrow$ JSON文件。
    \item \bff{GUI状态}：Fastbot $\rightarrow$ XML字符串 $\rightarrow$ \texttt{etree.Element} $\rightarrow$ XPath查询 $\rightarrow$ 布尔值 $\rightarrow$ 前置条件判断。
    \item \bff{执行结果}：\texttt{assert} 失败 $\rightarrow$ \texttt{AssertionError} $\rightarrow$ \texttt{result.addFailure()} $\rightarrow$ \texttt{PropertyExecutionInfo} $\rightarrow$ JSON文件 $\rightarrow$ HTML报告。
\end{itemize}

\subsection{关键函数/方法调用关系详细说明}

\subsubsection{性质检查核心调用链}

\begin{Verbatim}[fontsize=\footnotesize]
KeaTestRunner.run()
  |- collectAllProperties(test)
  |   `- iter_tests(suite)  # 递归遍历TestSuite
  |       `- hasattr(testMethod, PRECONDITIONS_MARKER)  # 检查@precondition标记
  |           `- self.allProperties[testMethodName] = t  # 存储性质
  |
  |- FastbotManager.start()  # 启动Fastbot服务（Android端）
  |- FastbotManager.init()   # 初始化配置
  |
  `- while循环（探索循环）:
      |- FastbotManager.stepMonkey()  # 发送Monkey事件
      |   `- HTTP POST /stepMonkey -> 返回XML hierarchy字符串
      |
      |- getValidProperties(xml_raw, result)
      |   |- U2Driver.getStaticChecker(hierarchy=xml_raw)
      |   |   `- U2StaticChecker.setHierarchy(hierarchy)
      |   |       |- etree.fromstring()  # 解析XML
      |   |       `- _HindenWidgetFilter()  # 过滤被覆盖控件
      |   |
      |   `- 遍历allProperties，执行precond(test)
      |       |- setattr(test, driverName, staticCheckerDriver)  # 依赖注入
      |       `- precond(test)  # 执行lambda函数
      |           `- self.d(...).exists
      |               `- StaticU2UiObject.exists
      |                   |- selector_to_xpath()  # 转换Selector为XPath
      |                   `- session.xml.xpath(xpath)  # 查询XML
      |
      |- 概率过滤 + 随机选择性质
      |
      |- U2Driver.getScriptDriver(mode="proxy")  # 获取可执行操作的driver
      |- setattr(test, driverName, scriptDriver)  # 依赖注入
      |
      |- test(result)  # 执行性质测试方法
      |   |- unittest.TestCase.__call__(result)
      |   |   |- 执行方法体（交互场景I）
      |   |   |   `- self.d.click() 等操作（U2ScriptDriver）
      |   |   |
      |   |   `- assert语句（后置条件Q）
      |   |       `- 如果失败 -> AssertionError
      |   |
      |   `- 捕获异常 -> result.addFailure/addError()
      |
      |- JsonResult.addFailure/addError()
      |   `- 更新PropertyExecutionInfo（state, tb）
      |
      `- result.flushResult()
          |- 写入result_*.json（统计信息）
          `- 写入property_exec_info_*.json（执行信息）
\end{Verbatim}

\subsubsection{数据流转路径详细说明}

\bff{性质对象流转}：
\begin{Verbatim}[fontsize=\footnotesize]
用户编写的性质文件 (properties/*.py)
  ↓ (unittest加载)
TestCase实例 (标记@precondition装饰器)
  ↓ (collectAllProperties收集)
allProperties字典 {propName: TestCase}
  ↓ (getValidProperties过滤)
validProps字典 {propName: TestCase}
  ↓ (概率过滤 + 随机选择)
execPropName, test (选中的性质)
  ↓ (执行)
test(result) 调用
  ↓ (记录)
PropertyExecutionInfo (startStepsCount, propName, state, tb)
  ↓ (持久化)
result.json + property_exec_info_*.json
  ↓ (报告生成)
BugReportGenerator解析 → HTML报告
\end{Verbatim}

\bff{GUI状态流转}：
\begin{Verbatim}[fontsize=\footnotesize]
Fastbot (Android端，通过Accessibility Service)
  ↓ (stepMonkey HTTP请求)
XML hierarchy字符串 (UI控件树)
  ↓ (U2StaticChecker.setHierarchy)
etree.Element (解析后的XML树)
  ↓ (_HindenWidgetFilter)
过滤后的XML树 (标记covered属性)
  ↓ (StaticU2UiObject.exists查询)
U2 Selector → XPath转换
  ↓ (xpath查询)
匹配的XML节点列表
  ↓ (bool转换)
exists属性值 (True/False)
  ↓ (precond lambda返回)
前置条件满足判断 (True/False)
\end{Verbatim}

\section{前沿技术关联}

\subsection{动态程序分析（Dynamic Program Analysis）}
Kea2的性质检查模块是动态程序分析的典型应用：
\begin{itemize}
    \item \bff{运行时监控}：在APP实际执行过程中收集程序状态（GUI层次结构、控件属性）。
    \item \bff{状态快照}：通过 \texttt{dumpHierarchy()} 获取执行点的完整GUI状态（XML hierarchy）。
    \item \bff{执行轨迹记录}：记录Monkey事件序列和性质执行信息，用于bug复现。
\end{itemize}

\subsection{运行时验证（Runtime Verification）}
该模块实现了轻量级的运行时验证框架：
\begin{itemize}
    \item \bff{性质规约语言}：PDL（基于Python装饰器）提供灵活的性质定义机制。
    \item \bff{在线监控}：在APP执行过程中实时检查性质，而非离线分析。
    \item \bff{违反检测}：通过断言评估自动检测性质违反，无需人工介入。
\end{itemize}

\subsection{程序插装（Program Instrumentation）}
虽然Kea2主要使用外部工具（Uiautomator2）获取状态，但其架构体现了插装思想：
\begin{itemize}
    \item \bff{代理模式}：通过Fastbot代理（port 8090）拦截和控制APP执行。
    \item \bff{事件注入}：Monkey事件可视为对APP的"插装"，驱动执行路径。
    \item \bff{状态提取}：通过Android Accessibility Service获取GUI状态，类似于插装点。
\end{itemize}

\subsection{基于性质的测试（Property-based Testing）}
这是该模块的核心技术范式：
\begin{itemize}
    \item \bff{性质驱动}：测试用例由性质规约（P$\rightarrow$I$\rightarrow$Q）驱动，而非固定输入。
    \item \bff{自动生成输入}：通过探索策略自动生成GUI事件序列。
    \item \bff{预言机制}：后置条件/断言作为测试预言，判断功能正确性。
\end{itemize}

\subsection{符号执行（部分关联）}
虽然Kea2主要使用具体执行（concrete execution），但其"主路径引导探索"策略体现了符号执行的启发：
\begin{itemize}
    \item \bff{路径推导}：基于性质的前置条件推导关键执行路径。
    \item \bff{目标导向}：优先探索可能触发目标性质的路径，而非盲目随机。
\end{itemize}

\subsection{混合测试（Hybrid Testing）}
Kea2支持传统单元测试与性质测试的混合执行：
\begin{itemize}
    \item \bff{HybridTestRunner}：允许在unittest中调用 \texttt{kea2\_breakpoint()} 切换到性质测试。
    \item \bff{无缝切换}：同一测试框架下支持两种测试范式，提升测试灵活性。
\end{itemize}

\section{总结}

Kea2的性质检查模块通过PDL语言定义性质、动态探索驱动执行、实时状态收集与断言评估，实现了Android应用功能bug的自动化检测。其核心创新在于将"用户指定性质"与"高效探索策略"结合，解决了传统工具"测试预言缺失"和"覆盖不足"的问题。

\bff{关键技术机制}：
\begin{enumerate}
    \item \bff{暂停执行与控制权转移}：当检测到满足前置条件的性质时，暂停Fastbot的探索循环，将控制权转移到性质测试方法的执行。
    \item \bff{依赖注入}：通过 \texttt{setattr} 动态注入不同的driver（静态检查器用于前置条件，脚本驱动用于性质执行）。
    \item \bff{状态快照}：通过XML hierarchy获取完整的GUI状态，支持高效的前置条件检查。
\end{enumerate}

该模块在97个历史bug中检测率达到94.8\%，并在实际应用中发现了25个新bug，验证了其有效性。

\end{document}